\documentclass{article}
\usepackage[utf8]{inputenc}
\usepackage[spanish]{babel}
\usepackage{amsmath}
\usepackage{graphicx}
\usepackage{geometry}
\usepackage{float}
\usepackage[american]{circuitikz}
\usetikzlibrary{babel}
\usepackage[hidelinks]{hyperref}
\usepackage{multirow}
\usepackage{array}
\usepackage{booktabs}
\usepackage{siunitx}
\usepackage{microtype}
\sisetup{
    output-decimal-marker = {.},
    group-separator = {,},
    group-minimum-digits = 4,
    inter-unit-product = \ensuremath{{}\cdot{}}
}

\geometry{a4paper, margin=1in}

\numberwithin{equation}{subsection}
\numberwithin{figure}{section}
\numberwithin{table}{section}

\begin{document}

\begin{center}
    \large\textbf{PRÁCTICA 2}
    \vspace{0.5cm}
    
    \huge\textbf{CIRCUITO PUENTE EQUILIBRADO}
\end{center}

\tableofcontents
\newpage

\section{OBJETIVOS}
\begin{enumerate}
    \item Hallar por análisis las relaciones matemáticas entre las resistencias en un circuito puente equilibrado.
    \item Aplicar experimentalmente el principio del puente equilibrado para medir una resistencia desconocida.
\end{enumerate}

\section{DESARROLLO TEÓRICO}

\subsection{Información Preliminar}

El volt-ohm-miliamperímetro (VOM) y el voltímetro electrónico (excepto el DVOM) son instrumentos de lectura directa 
en los cuales se utiliza un galvanómetro.

Las cantidades medidas, voltios, miliamperios u ohmios se leen directamente en una escala calibrada en la que se 
desplaza la aguja.

Para medir resistencias con más exactitud que la posible con un óhmetro se utiliza un puente. El puente de 
resistencias tiene un galvanómetro de alta sensibilidad como dispositivo indicador, una resistencia variable 
calibrada que sirve de patrón de referencia y una fuente de tensión en una disposición adecuada de circuito. 
El galvanómetro sirve como indicador de cero y pone de manifiesto la condición de equilibrio. La resistencia 
se lee en una escala calibrada asociada a la resistencia variable y a un multiplicador de escala o margen.

\begin{figure}[!htbp]
    \centering
    \centering
    \begin{circuitikz}[scale=1.2, american]
        % Define styles
        \ctikzset{bipoles/length=1.2cm}
        
        % Nodes
        \coordinate (D) at (0,0);
        \coordinate (C) at (4,2);
        \coordinate (B) at (0,4);
        \coordinate (A) at (-4,2);

        % Bridge Arms
        \draw (B) to[R, l_=$R_1$, i>=$I_1$, *-*] (A);
        \draw (B) to[vR, l=$R_3$, i=$I_3$, *-*] (C); % Variable (Brazo patrón)
        \draw (A) to[R, l_=$R_2$, i>=$I_2$, *-*] (D);
        \draw (C) to[R, l=$R_x$, i=$I_x$, *-*] (D);

        % Galvanometer
        \draw (A) to[rmeter, t=G, *-*] (C);

        % Power Supply (External)
        \draw (B) -- (0,5) -- (5,5) -- (5,2) 
              to[V, l=V, invert] (5,0) -- (5,-1) -- (0,-1) -- (D);

        % Labels
        \node[left] at (A) {A};
        \node[above] at (B) {B};
        \node[right] at (C) {C};
        \node[below] at (D) {D};
    \end{circuitikz}
    \caption{Puente de Wheatstone para Medir la Resistencia Desconocida $R_x$.}
    \label{fig:25-1}
\end{figure}

La figura~\ref{fig:25-1} representa el tipo más corriente, el puente de Wheatstone. En esta disposición 
hay cuatro resistencias en el circuito. La resistencia desconocida que se desea medir, $R_x$, se conecta 
entre los bornes C y D. $R_1$ y $R_2$ son resistencias de precisión de valor fijo llamadas «brazos de 
proporción», y $R_3$ es una resistencia variable conocida. El indicador es un galvanómetro de alta 
sensibilidad con el cero en el centro, G. La corriente pasará por G cuando haya una diferencia de 
potencial entre los puntos A y C. Cuando no existe diferencia de potencial entre estos puntos, es decir, 
cuando $V_{AC} = 0$, la aguja del galvanómetro volverá a su posición de reposo o cero. Cuando $V_{AC} = 0$, 
se dice que el puente está equilibrado.

En el circuito de la figura~\ref{fig:25-1} sean $I_1$ la corriente en $R_1$, $I_2$ en $R_2$, $I_3$ en 
$R_3$ e $I_x$ en $R_x$. La tensión de A a B es $I_1 \times R_1$, es decir

\begin{align}
    V_{AB} &= I_1 R_1 \nonumber \\
    \text{Análogamente } V_{AD} &= I_2 R_2 \label{eq:25-1} \\
    V_{CB} &= I_3 R_3 \label{eq:25-2} \\
    V_{CD} &= I_x R_x \label{eq:25-3}
\end{align}

Para obtener el equilibrio ($V_{AC} = 0$), la diferencia de tensión entre A y B debe ser igual a la existente 
entre C y B. O sea,
\begin{equation}
    V_{AB} = V_{CB} \label{eq:25-4}
\end{equation}

Por tanto, $I_1 R_1 = I_3 R_3$, y
\begin{equation}
    \frac{I_1}{I_3} = \frac{R_3}{R_1} \label{eq:25-5}
\end{equation}

Análogamente $V_{AD} = V_{CD}$
\begin{equation}
    \frac{I_2}{I_x} = \frac{R_x}{R_2} \label{eq:25-6}
\end{equation}

En el equilibrio no pasa corriente por G, por tanto
\begin{align}
    I_1 &= I_2 \label{eq:25-7} \\
    I_3 &= I_x \label{eq:25-8}
\end{align}

Sustituyendo los resultados de las ecuaciones \eqref{eq:25-7} y \eqref{eq:25-8} en las \eqref{eq:25-5} y \eqref{eq:25-6}, hallaremos
\begin{equation}
    \frac{I_1}{I_x} = \frac{R_3}{R_1} \label{eq:25-9}
\end{equation}
\begin{equation}
    \frac{I_1}{I_x} = \frac{R_x}{R_2} \label{eq:25-10}
\end{equation}

por tanto
\begin{equation}
    \frac{R_x}{R_2} = \frac{R_3}{R_1} \label{eq:25-11}
\end{equation}
y
\begin{equation}
    R_x = \left( \frac{R_2}{R_1} \right) \times R_3 \label{eq:25-12}
\end{equation}

La ecuación fundamental \eqref{eq:25-12} enuncia que con el puente equilibrado la resistencia desconocida $R_x$ 
es igual al producto de la razón o cociente de los brazos de proporción ($R_2/R_1$) y el brazo de 
comparación $R_3$.

Se obtienen la máxima exactitud y la máxima sensibilidad cuando $R_1 = R_2$ es decir, cuando la razón 
$R_2/R_1 = 1$. En esta condición $R_x = R_3$. Si $R_3$ es un reóstato del tipo de décadas de alta precisión, 
el valor $R_x$ se puede leer directamente en la escala calibrada del reóstato cuando se ajusta $R_3$ para 
el equilibrio.

\subsection{RESUMEN}
\begin{enumerate}
    \item Para medir valores de resistencia con mayor exactitud que la obtenible con un volt-ohm-miliamperímetro 
    o un voltímetro electrónico se utiliza un puente de Wheatstone.
    \item Los brazos de proporción $R_1$ y $R_2$ de la figura~\ref{fig:25-1} son resistencias de precisión que se 
    conmutan para que den el multiplicador deseado del margen de resistencia del puente.
    \item El brazo patrón o de comparación $R_3$ es una resistencia de precisión calibrada y variable por medio de 
    la cual se equilibra el puente cuando se mide una resistencia desconocida $R_x$. El valor de $R_x$ se lee, en 
    el equilibrio, en la escala calibrada del brazo patrón, y se multiplica por la relación de los brazos 
    de proporción.
    \item Como indicador se utiliza un galvanómetro de alta sensibilidad con 0 en el centro. El puente se alimenta 
    con una fuente de tensión.
    \item Cuando la tensión entre los terminales del galvanómetro ($V_{AC}$ en la tabla~\ref{tab:25-1}) es cero, no 
    hay corriente en el galvanómetro, y por lo tanto éste dará lectura cero. Es en esta posición de equilibrio en 
    la que se lee el valor de la resistencia desconocida.
    \item Estando en equilibrio un puente de Wheatstone, los productos de las resistencias de los brazos opuestos son 
    iguales, es decir, en la tabla~\ref{tab:25-1}
    \begin{equation} R_x \times R_1 = R_2 \times R_3 \label{eq:resumen_1} \end{equation}
    y en esta condición el valor de $R_x$ se deduce de
    \begin{equation} R_x = \frac{R_2 \times R_3}{R_1} \label{eq:resumen_2} \end{equation}
    \item Por consiguiente, $R_x$ es igual al producto de la relación de los brazos de proporción por la resistencia 
    del brazo patrón.
\end{enumerate}

\newpage

\section{MATERIALES NECESARIOS}
Fuente de alimentación: Fuente de c.c. de tensión regulada variable.

\textbf{Equipo:} Galvanómetro con cero en el centro o un VOM de \SI{20000}{\ohm\per\volt}; voltímetro electrónico.

\textbf{Resistencias:} 1/2 W, dos de \SI{5100}{\ohm} (\SI{5}{\percent}); surtido de seis valores diferentes de resistencias 
(valor mínimo 500, valor máximo menos que \SI{10000}{\ohm}).

\textbf{Opcional:} Otras resistencias como las necesarias en la operación 5.

\textbf{Diversos:} Cajas de décadas de resistencia, variable desde 1 a 10.000 o potenciómetro de \SI{10000}{\ohm}; dos 
interruptores unipolares.

\begin{figure}[!htbp]
    \centering
    \centering
    \begin{circuitikz}[scale=1.2, american]
        % Coordinates
        \coordinate (D) at (0,0);
        \coordinate (C) at (4,2);
        \coordinate (B) at (0,4);
        \coordinate (A) at (-4,2);

        % Bridge
        \draw (B) to[R, l_=$R_1$, *-*] (A);
        \draw (B) to[vR, l=$R_3$, *-*] (C);
        \draw (A) to[R, l_=$R_2$, *-*] (D);
        \draw (C) to[R, l=$R_x$, *-*] (D);

        % Central Branch with Galvanometer shunting
        % Creating a more readable schematic for the shunt
        \draw (A) -- (-1.5, 2); 
        \draw (-1.5, 2) to[rmeter, t=M, *-*] (1.5, 2); % Galvanometer
        % Shunt path
        \draw (-1.5, 2) -- (-1.5, 2.8) to[R, l=$R_4$] (1.5, 2.8) -- (1.5, 2);
        \draw (-1.5, 2) -- (-1.5, 1.2) to[nos, l_=$S_2$] (1.5, 1.2) -- (1.5, 2);
        
        \draw (1.5, 2) -- (C);

        % Power Supply
        \draw (B) -- (0,5) -- (-5,5) to[V, l_=6V] (-5,-1) 
              to[nos, l=$S_1$] (0,-1) -- (D);

        % Labels
        \node[left] at (A) {A};
        \node[above] at (B) {B};
        \node[right] at (C) {C};
        \node[below] at (D) {D};
    \end{circuitikz}
    \caption{Puente de Wheatstone Experimental.}
    \label{fig:25-2}
\end{figure}

\section{PROCEDIMIENTO}
\begin{enumerate}
    \item Conectar el circuito (figura~\ref{fig:25-2}). $S_1$ y $S_2$ están abiertos. Ajustar la alimentación regulada para 
    V = \SI{6}{\volt}. Ajustar $R_3$, caja de décadas de resistencia, para resistencia máxima. (Si no se dispone de 
    caja de década, utilizar un potenciómetro de \SI{10000}{\ohm}.) M es un galvanómetro con cero en el centro. 
    (Si no se dispone de él, utilizar el margen o alcance de corriente de más sensibilidad en el VOM de \SI{20000}{\ohm\per\volt}, 
    es decir, alcances de \SI{50}{\micro\ampere} ó \SI{60}{\micro\ampere}.) El objeto de $R_4$ es reducir la sensibilidad 
    de M durante las fases iniciales de ajuste.
    \item Conectar $R_x$, resistencia a medir, entre C y D. ($R_x$ no debe ser menor de \SI{560}{\ohm}, ni mayor de \SI{10000}{\ohm}.) 
    Cerrar $S_1$. $S_2$ está aún abierto. La aguja del instrumento se desviará hacia arriba. De lo contrario, mirar si 
    el brazo de referencia o comparación $R_3$ está en el máximo.
    \item Reducir la resistencia de $R_3$ gradualmente hasta que la lectura de M se aproxime a cero. Cerrar $S_2$ y 
    reajustar $R_3$ si es necesario hasta que M dé lectura cero. Leer el valor de $R_x$ en la caja de décadas de 
    resistencia calibrada y anotarla en la tabla~\ref{tab:25-1}. Anotar también el valor de $R_x$ indicado por el código de color 
    y su tolerancia nominal.
    
    \textbf{NOTA:} Si $R_3$ es un potenciómetro de \SI{10000}{\ohm} conectado como reóstato en lugar de una caja de 
    décadas, eliminar $R_3$ en el circuito sin variar su ajuste y medirla con un óhmetro. Anotar esto como valor 
    de $R_x$ en la tabla~\ref{tab:25-1}.
    
    \item Abrir $S_1$ y $S_2$. Ajustar nuevamente $R_3$ para resistencia máxima. Quitar la primera $R_x$ del circuito. 
    Por el procedimiento de circuito puente de las operaciones 2 y 3, medir sucesivamente cada una de las cuatro 
    resistencias recibidas. Anotar los valores medidos y codificados en la tabla~\ref{tab:25-1}.
\end{enumerate}

\begin{table}[!htbp]
    \centering
    \caption{Medición en el Puente de $R_x$; $R_2/R_1 = 1$}
    \label{tab:25-1}
    \begin{tabular}{lcccccc}
        \toprule
        \textbf{Resistencia n.º} & \textbf{1} & \textbf{2} & \textbf{3} & \textbf{4} & \textbf{5} & \textbf{6} \\ \midrule
        Valor nominal (código) & & & & & & \\ \midrule
        Por ciento de tolerancia & & & & & & \\ \midrule
        Valor medido & & & & & & \\ \bottomrule
    \end{tabular}
\end{table}

\subsection{Adicional — Problema de diseño}
\begin{enumerate}
    \setcounter{enumi}{4}
    \item Utilizando la misma $R_3$ como brazo de comparación, proyectar un circuito puente con el que se puedan medir 
    resistencias de hasta
    \begin{enumerate}
        \item \SI{50000}{\ohm} (aproximadamente)
        \item \SI{250000}{\ohm} (aproximadamente).
    \end{enumerate}
    Dibujar los esquemas correspondientes a a y b indicando los valores necesarios de todos los componentes 
    del puente. Tomar los componentes necesarios. Si no se dispone de resistencias de \SI{5}{\percent} de tolerancia, 
    medir las resistencias de los brazos de proporción con un óhmetro y calcular la relación de multiplicación 
    por los valores medidos. Con el puente a medir tres resistencias mayores de \SI{10000}{\ohm}, pero menores 
    de \SI{50000}{\ohm}. Con el puente b medir tres resistencias mayores de \SI{50000}{\ohm} y menores de \SI{250000}{\ohm}.
    Anotar los resultados de las mediciones en tablas separadas análogas a la tabla~\ref{tab:25-1}. Identificar la relación 
    de multiplicador en cada caso.
\end{enumerate}

\section{PREGUNTAS}
\begin{enumerate}
    \item ¿Qué es lo que determina la aproximación de las mediciones hechas con el puente experimental de resistencias?
    \item ¿Por qué es necesario emplear un galvanómetro muy sensible como indicador en esta práctica?
    \item ¿Cuál es la finalidad de la fuente de tensión V? ¿Es crítico el nivel de tensión (V = \SI{6}{\volt}) de la 
    fuente de tensión?
    \item ¿Corresponden los resultados de sus mediciones a los valores nominales de las resistencias medidas? De lo 
    contrario, explicar el motivo.
\end{enumerate}

\subsection{Adicional}
\begin{enumerate}
    \setcounter{enumi}{4}
    \item En el equilibrio, $R_1 = R_2 = \SI{5000}{\ohm}$; $R_x = R_3 = \SI{6000}{\ohm}$; $R_m = \SI{3000}{\ohm}$. Hallar la 
    resistencia total conectada a la fuente de alimentación (V). Presentar los cálculos.
\end{enumerate}

\end{document}